\setcounter{section}{-1}
\section{\S 3.0 基础知识复习}\label{sect:3.0基础知识复习}

    \subsection{电解质与非电解质} \label{subsec:电解质与非电解质}

        在水溶液\uline{或}熔融状态下\(\overset{\text{\textcolor{blue}{电离}}}{\text{导电}}\)的\uline{化合物}叫做电解质；
        
        在水溶液\uline{和}熔融状态下都不能\(\overset{\text{\textcolor{blue}{电离}}}{\text{导电}}\)的\uline{化合物}叫做非电解质。

        \textcolor{red}{必须是化合物本身电离，而非与水反应后产生离子}

        \textcolor{red}{是物质具有的性质，与物质所处状态无关}

        \textcolor{blue}{{\bfseries 电解质}：酸、碱、盐、水}

        \textcolor{blue}{电解质导电原因：有\uline{自由移动}的离子}

        \qquad 电离：电解质产生自由移动离子的过程

        \textcolor{red}{\(\therefore\text{溶液的导电能力}\propto\text{离子浓度}\)}

        离子浓度\(\left\{\begin{matrix}
            \text{溶液浓度}\\
            \text{电离程度}
        \end{matrix}\right.\)

    \subsection{强弱电解质}\label{subsec:强弱电解质}

        强电解质——在水溶液中\textcolor{blue}{完全电离}为离子的电解质
        
        \phantom{强电解质——}电离方程式用``\cemh{\xleq{}}''\quad 只存在离子

        弱电解质——在水溶液中\textcolor{blue}{部分电离}为离子的电解质
        
        \phantom{强电解质——}电离方程式用``\cemh{<=>}''\quad 离子、分子都有

        强电解质：强酸、强碱、绝大多数盐

        弱电解质：弱酸、弱碱、水

        \begin{itemize}
            \item 常见强酸：\cemh{H2SO4}、\cemh{HNO3}、\cemh{HCl}、\cemh{HBr}、\cemh{HI}、\cemh{HClO4}
            \item 常见强碱：\cemh{KOH}、\cemh{NaOH}、\cemh{Ba(OH)2}、\cemh{Ca(OH)2}
            \item 常见弱酸：见 \textbf{表~\ref{tab:弱酸-课本3.3}}
            \item 常见弱碱：氨水、不溶性碱
        \end{itemize}

        \begin{table}[htbp]
            \centering
            \caption{常见弱酸、弱碱的电离平衡常数（\SI{25}{\degreeCelsius}）}
            \label{tab:弱酸-课本3.3}
            \begin{tabular}{|c|c||c|c|}
                \toprule
                \small 一元弱酸{\footnotesize 或}弱碱 &\small \( K_\text{a} \) {\footnotesize 或} \( K_\text{b} \) &\small 多元弱酸 &\small \( K_\text{a1} \) {\footnotesize 和} \( K_\text{a2} \) \\
                \midrule\footnotesize
                \cemh{HCOOH}\footnotemark &\footnotesize \SI{1.8E-4}{} &\footnotesize 
                \multirow{2}[0]{*}{\cemh{H2CO3}} &\footnotesize \( K_\text{a1} = \SI{4.2E-7}{} \) \\\footnotesize
                \cemh{HNO2} &\footnotesize \SI{7.2E-4}{} &\footnotesize 
                                                    &\footnotesize \( K_\text{a2} = \SI{4.8E-11}{} \) \\\footnotesize 
                \cemh{HF} &\footnotesize \SI{5.6E-4}{} &\footnotesize 
                \multirow{2}[0]{*}{\cemh{H2SO3}} &\footnotesize \( K_\text{a1} = \SI{1.2E-2}{} \) \\\footnotesize
                \cemh{CH3COOH} &\footnotesize \SI{1.8E-5}{} &\footnotesize 
                                                    &\footnotesize \( K_\text{a2} = \SI{6.2E-8}{} \) \\\footnotesize
                \cemh{HCN}\footnotemark &\footnotesize \SI{4.0E-10}{} &\footnotesize 
                \multirow{2}[0]{*}{\chemfig{COOH(-[2, .7]COOH)}}\footnotemark &\footnotesize \( K_\text{a1} = \SI{5.6E-2}{} \) \\\footnotesize 
                \cemh{NH3 . H2O} &\footnotesize \SI{1.8E-5}{} &\footnotesize 
                                                                                &\footnotesize \( K_\text{a2} = \SI{5.4E-5}{} \) \\
                \midrule\footnotesize
                \multirow{2}[0]{*}{\cemh{H2S}} &\footnotesize \( K_\text{a1} = \SI{5.7E-8}{} \) &\footnotesize 
                \multirow{3}[0]{*}{\cemh{H3PO4}} &\footnotesize \( K_\text{a1} = \SI{7.52E-3}{} \)  \\\footnotesize
                                                &\footnotesize \( K_\text{a2} = \SI{1.2E-15}{} \) &\footnotesize 
                                                    &\footnotesize \( K_\text{a2} = \SI{6.23E-8}{} \) \\\footnotesize
                                                &&&\footnotesize \( K_\text{a3} = \SI{2.2E-13}{} \) \\
                \bottomrule
            \end{tabular}
        \end{table}
        \footnotetext[1]{甲酸}\footnotetext[2]{氢氰酸}\footnotetext[3]{草酸}
    
    \subsection{电解质与离子、共价化合物的关系}\label{subsec:电解质与离子、共价化合物的关系}

        离子化合物——强电解质

        共价化合物\(\left\{\begin{matrix}
            \text{强极性键\hspace{1em}} & \text{——强酸——} & \text{强电解质}\\
            \text{较强极性键} & \text{——\phantom{强酸}——} & \text{弱电解质}\\
            \text{弱极性键\hspace{1em}} & \text{——\phantom{强酸}——} & \text{非电解质}
        \end{matrix}\right.\)

        离子化合物——水溶液、熔融状态电离

        共价化合物——熔融状态电离
    